% answers.tex —— 1999 年·上海卷·理工农医类 · 参考答案与详细解析片段
% 由 tools/generate_projects.py 依来源总账自 1999/1999_数学真题_上海理_试题.md + 1999/1999_数学真题_上海理_解析正文.md + 1999/1999_数学真题_上海理_答案要点.md 生成，勿手改（改源 MD 后重跑）


\begin{answer}
\gksolhead{1}{}

\textbf{题目}：

$\operatorname{tg}\left(\arccos\frac{\sqrt2}{2}-\frac\pi6\right)=$ \gkblank{}。

\Needspace{4\baselineskip}
\textbf{解析}

\Needspace{5\baselineskip}
$\arccos(\sqrt2/2)=\pi/4$，故原式为

\nopagebreak[4]
\[
\tan\left(\frac\pi4-\frac\pi6\right)=\tan\frac\pi{12}=2-\sqrt3.
\]

\gkanswer{答案：$2-\sqrt3$。}

\end{answer}
\begin{answer}
\gksolhead{2}{}

\textbf{题目}：

函数 $f(x)=\log_2 x+1\quad (x\ge4)$ 的反函数 $f^{-1}(x)$ 的定义域是 \gkblank{}。

\Needspace{4\baselineskip}
\textbf{解析}

原函数只在 $x\ge4$ 上取值，且单调递增；其值域为 $[\log_2 4+1,+\infty)=[3,+\infty)$，即反函数的定义域。

\gkanswer{答案：$[3,+\infty)$。}

\end{answer}
\begin{answer}
\gksolhead{3}{}

\textbf{题目}：

在 $\left(x^3+\frac2{x^2}\right)^5$ 的展开式中，含 $x^5$ 项的系数为 \gkblank{}。

\Needspace{4\baselineskip}
\textbf{解析}

\Needspace{5\baselineskip}
取 $k$ 个因子 $2x^{-2}$，二项式通项的 $x$ 次幂为 $3(5-k)-2k=15-5k$。令 $15-5k=5$ 得 $k=2$，对应系数

\nopagebreak[4]
\[
\binom52 2^2=40.
\]

\gkanswer{答案：$40$。}

\end{answer}
\begin{answer}
\gksolhead{4}{}

\textbf{题目}：

在 $\triangle ABC$ 中，若 $\angle B=30^\circ$，$AB=2\sqrt3$，$AC=2$，则 $\triangle ABC$ 的面积 $S$ 是 \gkblank{}。

\Needspace{4\baselineskip}
\textbf{解析}

\Needspace{5\baselineskip}
记 $b=AC=2$、$c=AB=2\sqrt3$、$B=30^\circ$。由正弦定理，$\sin C=c\sin B/b=\sqrt3/2$，所以 $C=60^\circ$ 或 $120^\circ$，相应 $A=90^\circ$ 或 $30^\circ$。于是

\nopagebreak[4]
\[
S=\frac12bc\sin A=2\sqrt3\quad\text{\text{或}}\quad\sqrt3.
\]

\gkanswer{答案：$2\sqrt3$ 或 $\sqrt3$。}

\end{answer}
\begin{answer}
\gksolhead{5}{}

\textbf{题目}：

若平移坐标系，将曲线方程 $y^2+4x-4y-4=0$ 化为标准方程，则坐标原点应移到点 $O'(\underline{\quad},\underline{\quad})$。

\Needspace{4\baselineskip}
\textbf{解析}

\Needspace{5\baselineskip}
配方得

\nopagebreak[4]
\[
y^2+4x-4y-4=0\iff (y-2)^2=-4(x-2).
\]

标准抛物线顶点是 $(2,2)$，即原点移到 $O'(2,2)$。

\gkanswer{答案：$(2,2)$。}

\end{answer}
\begin{answer}
\gksolhead{6}{}

\textbf{题目}：

函数 $y=2\sin x\cos x-2\sin^2x+1$ 的最小正周期 $T=$ \gkblank{}。

\Needspace{4\baselineskip}
\textbf{解析}

利用 $2\sin x\cos x=\sin2x$ 与 $1-2\sin^2x=\cos2x$，原函数化为 $y=\sin2x+\cos2x=\sqrt2\sin(2x+\pi/4)$，最小正周期为 $\pi$。

\gkanswer{答案：$\pi$。}

\end{answer}
\begin{answer}
\gksolhead{7}{}

\textbf{题目}：

函数 $y=2\sin\left(2x+\frac\pi6\right)\quad (x\in[-\pi,0])$ 的单调递减区间是 \gkblank{}。

\Needspace{4\baselineskip}
\textbf{解析}

\Needspace{5\baselineskip}
导数 $y'=4\cos(2x+\pi/6)$。导数在区间内部为负、两端为零；与给定定义域 $[-\pi,0]$ 相交后，函数在区间

\nopagebreak[4]
\[
\left[-\frac{5\pi}{6},-\frac\pi3\right].
\]

答案要点按卷内记作开区间；端点处导数虽为零，函数在闭区间上仍严格递减，两种记法等价。

\gkanswer{答案：$\left(-\dfrac{5\pi}6,-\dfrac{\pi}3\right)$（源作开区间；函数在闭区间上同样严格递减）。}

\end{answer}
\begin{answer}
\gksolhead{8}{}

\textbf{题目}：

若将向量 $\vec a=\{2,1\}$ 围绕原点按逆时针方向旋转 $\frac\pi4$ 得到向量 $\vec b$，则向量 $\vec b$ 的坐标为 \gkblank{}。

\Needspace{4\baselineskip}
\textbf{解析}

\Needspace{5\baselineskip}
将 $(2,1)$ 逆时针旋转 $\pi/4$：

\nopagebreak[4]
\[
\begin{pmatrix}\cos\frac\pi4&-\sin\frac\pi4\\\sin\frac\pi4&\cos\frac\pi4\end{pmatrix}
\begin{pmatrix}2\\1\end{pmatrix}
=\begin{pmatrix}\frac{\sqrt2}{2}\\\frac{3\sqrt2}{2}\end{pmatrix}.
\]

按原题坐标花括号记法，答案：$\{\sqrt2/2,3\sqrt2/2\}$。

\gkanswer{答案：$\left\{\dfrac{\sqrt2}{2},\dfrac{3\sqrt2}{2}\right\}$。}

\end{answer}
\begin{answer}
\gksolhead{9}{}

\textbf{题目}：

极坐标方程 $5\rho^2\cos2\theta+\rho^2-24=0$ 所表示的曲线焦点的极坐标为 \gkblank{}。

\Needspace{4\baselineskip}
\textbf{解析}

\Needspace{5\baselineskip}
用 $\rho^2=x^2+y^2$、$\rho^2\cos2\theta=x^2-y^2$ 化为

\nopagebreak[4]
\[
6x^2-4y^2=24\iff\frac{x^2}{4}-\frac{y^2}{6}=1.
\]

故 $a^2=4$、$b^2=6$、$c^2=a^2+b^2=10$，焦点为 $(\pm\sqrt{10},0)$。转成原卷采用的极坐标表示为 $(\sqrt{10},0)$、$(\sqrt{10},\pi)$。

\gkanswer{答案：$(\sqrt{10},0)$、$(\sqrt{10},\pi)$。}

\end{answer}
\begin{answer}
\gksolhead{10}{}

\textbf{题目}：

在等差数列 $\{a_n\}$ 中，满足 $3a_4=7a_7$，且 $a_1>0$，$S_n$ 是数列 $\{a_n\}$ 前 $n$ 项的和。若 $S_n$ 取得最大值，则 $n=$ \gkblank{}。

\Needspace{4\baselineskip}
\textbf{解析}

\Needspace{5\baselineskip}
设公差为 $d$。由 $3(a_1+3d)=7(a_1+6d)$ 得 $d=-4a_1/33<0$。又

\nopagebreak[4]
\[
a_n=a_1+(n-1)d=a_1\left(1-\frac{4(n-1)}{33}\right),
\]

所以 $a_9=a_1/33>0$，而 $a_{10}=-a_1/11<0$。部分和在第9项仍增加、从第10项起减少，最大值取在 $n=9$。

\gkanswer{答案：$9$。}

\end{answer}
\begin{answer}
\gksolhead{11}{}

\textbf{题目}：

若连续掷两次骰子得到的点数 $m,n$ 是点 $P$ 的坐标，则点 $P$ 落在圆 $x^2+y^2=16$ 内的概率是 \gkblank{}。

\Needspace{4\baselineskip}
\textbf{解析}

样本空间是 $6\times6=36$ 个等可能有序点。按 $m=1,2,3,4,5,6$ 分组，满足 $m^2+n^2<16$ 的 $n$ 数依次为 $3,3,2,0,0,0$，共8个。

\gkanswer{答案：$\dfrac29$。}

\end{answer}
\begin{answer}
\gksolhead{12}{}

\textbf{题目}：

若四面体各棱的长是 $1$ 或 $2$，且该四面体不是正四面体，则其体积的值是 \gkblank{}（只需写出一个可能的值）。

\Needspace{4\baselineskip}
\textbf{解析}

\Needspace{5\baselineskip}
构造底面为边长1的等边三角形，取其中心为 $z$ 轴原点；底面三个顶点到中心的距离为 $1/\sqrt3$。令顶点 $D$ 位于中心正上方，且高

\nopagebreak[4]
\[
h=\sqrt{2^2-(1/\sqrt3)^2}=\sqrt{11/3}.
\]

\Needspace{5\baselineskip}
则三条侧棱均为2、底面三边均为1，构成非正四面体。其体积

\nopagebreak[4]
\[
V=\frac13\cdot\frac{\sqrt3}{4}\cdot\sqrt{\frac{11}{3}}=\frac{\sqrt{11}}{12}.
\]

\gkanswer{答案：$\dfrac{\sqrt{11}}{12}$、$\dfrac{\sqrt{14}}{12}$、$\dfrac{\sqrt{11}}6$（三者任一）。}

\end{answer}
\begin{answer}
\gksolhead{13}{}

\textbf{题目}：

直线 $y=\frac{\sqrt3}{3}x$ 绕原点按逆时针方向旋转 $30^\circ$ 后所得直线与圆 $(x-2)^2+y^2=3$ 的位置关系是\mbox{（　　）}

（A）直线过圆心；

（B）直线与圆相交，但不过圆心；

（C）直线与圆相切；

（D）直线与圆没有公共点。

\Needspace{4\baselineskip}
\textbf{解析}

原直线倾角为 $30^\circ$，逆时针转 $30^\circ$ 后斜率为 $\tan60^\circ=\sqrt3$。圆心 $(2,0)$ 到 $y=\sqrt3x$ 的距离为 $2\sqrt3/2=\sqrt3$，等于半径，故相切。

\gkanswer{答案：C。}

\end{answer}
\begin{answer}
\gksolhead{14}{}

\textbf{题目}：

下列以 $t$ 为参数的参数方程所表示的曲线中，与方程 $xy=1$ 所表示的曲线完全一致的是\mbox{（　　）}

（A）$\begin{cases}x=t^{1/2},\\ y=t^{-1/2};\end{cases}$
（B）$\begin{cases}x=|t|,\\ y=\frac1{|t|};\end{cases}$
（C）$\begin{cases}x=\cos t,\\ y=\sec t;\end{cases}$
（D）$\begin{cases}x=\operatorname{tg}t,\\ y=\operatorname{ctg}t.\end{cases}$

\Needspace{4\baselineskip}
\textbf{解析}

参数式D满足 $x=\tan t\ne0$、$y=\cot t=1/x$。对任意满足 $xy=1$ 的实数点，取 $t=\arctan x$ 即可得到该点，因此D覆盖完整双曲线。A、B只覆盖正支；C中 $x=\cos t$ 限于 $0<|x|\le1$，没有覆盖全曲线。

\gkanswer{答案：D。}

\end{answer}
\begin{answer}
\gksolhead{15}{}

\textbf{题目}：

若 $a<b<0$，则下列结论中正确的是\mbox{（　　）}

（A）不等式 $\frac1a>\frac1b$ 和 $\frac1{|a|}>\frac1{|b|}$ 均不能成立；

（B）不等式 $\frac1{a-b}>\frac1a$ 和 $\frac1{|a|}>\frac1{|b|}$ 均不能成立；

（C）不等式 $\frac1{a-b}>\frac1a$ 和 $\left(a+\frac1b\right)^2>\left(b+\frac1a\right)^2$ 均不能成立；

（D）不等式 $\frac1{|a|}>\frac1{|b|}$ 和 $\left(a+\frac1b\right)^2>\left(b+\frac1a\right)^2$ 均不能成立。

\Needspace{4\baselineskip}
\textbf{解析}

\Needspace{5\baselineskip}
因 $a<b<0$，有 $a-b>a$ 且 $a-b,a$ 均为负数，故 $1/(a-b)<1/a$；同时 $|a|>|b|$，故 $1/|a|<1/|b|$。所以选项B所列两个“大于”不等式都不可能成立。另将两平方项作差：

\nopagebreak[4]
\[
\left(a+\frac1b\right)^2-\left(b+\frac1a\right)^2
=(a-b)(a+b)\left(1+\frac1{ab}\right)^2>0,
\]

因为 $a-b<0$、$a+b<0$ 且 $ab>0$，该平方不等式对所有 $a<b<0$ 恒成立。故选项C、D中“均不能成立”的判断为假；选项A也为假，因为 $1/a>1/b$ 恒成立。正确答案为B。

\gkanswer{答案：B。}

\end{answer}
\begin{answer}
\gksolhead{16}{}

\textbf{题目}：

设有直线 $m,n$ 和平面 $\alpha,\beta$，则在下列命题中，正确的是\mbox{（　　）}

（A）若 $m\parallel n$，$m\subset\alpha$，$n\subset\beta$，则 $\alpha\parallel\beta$；

（B）若 $m\perp\alpha$，$m\perp n$，$n\subset\beta$，则 $\alpha\parallel\beta$；

（C）若 $m\parallel n$，$n\perp\beta$，$m\subset\alpha$，则 $\alpha\perp\beta$；

（D）若 $m\parallel n$，$m\perp\alpha$，$n\perp\beta$，则 $\alpha\perp\beta$。

\Needspace{4\baselineskip}
\textbf{解析}

若 $m\parallel n$ 且 $n\perp\beta$，则 $m\perp\beta$。又 $m\subset\alpha$，平面 $\alpha$ 含有一条垂直于平面 $\beta$ 的直线，故 $\alpha\perp\beta$，C正确。A、B、D均不能由题设推出所称平行/垂直关系。

\gkanswer{答案：C。}

\end{answer}
\begin{answer}
\gksolhead{17}{}

\textbf{题目}（本题满分10分）：

设集合 $A=\{x\mid |x-a|<2\}$，$B=\left\{x\mid\frac{2x-1}{x+2}<1\right\}$。若 $A\subseteq B$，求实数 $a$ 的取值范围。

\Needspace{4\baselineskip}
\textbf{解析}

\Needspace{5\baselineskip}
解不等式时需保留分母限制 $x\ne-2$：

\nopagebreak[4]
\[
\frac{2x-1}{x+2}<1\iff\frac{x-3}{x+2}<0\iff -2<x<3.
\]

故 $B=(-2,3)$，而 $A=(a-2,a+2)$。两个区间均为开区间，因此 $A\subseteq B$ 当且仅当 $a-2\ge-2$ 且 $a+2\le3$，即 $0\le a\le1$。

\gkanswer{答案：$0\le a\le1$。}

\end{answer}
\begin{answer}
\gksolhead{18}{}

\textbf{题目}（本题满分12分）：

\Needspace{5\baselineskip}
设正数数列 $\{a_n\}$ 为一等比数列，且 $a_2=4$，$a_4=16$，求
\nopagebreak[4]
\[
\lim_{n\to\infty}\frac{\lg a_{n+1}+\lg a_{n+2}+\cdots+\lg a_{2n}}{n^2}
\]

\Needspace{4\baselineskip}
\textbf{解析}

\Needspace{5\baselineskip}
等比数列各项为正，且 $q^2=a_4/a_2=4$，故公比 $q=2$；由 $a_2=4$ 得 $a_k=2^k$。于是

\nopagebreak[4]
\[
\frac{\lg a_{n+1}+\cdots+\lg a_{2n}}{n^2}
=\lg2\,\frac{(n+1)+\cdots+2n}{n^2}
=\lg2\,\frac{n(3n+1)}{2n^2}\longrightarrow\frac32\lg2.
\]

\gkanswer{答案：$\dfrac32\lg2$。}

\end{answer}
\begin{answer}
\gksolhead{19}{}

\textbf{题目}（本题满分14分）：

已知方程 $x^2+(4+i)x+4+ai=0\ (a\in\mathbb R)$ 有实根 $b$，且 $z=a+bi$，求复数 $\bar z(1-ci)\ (c>0)$ 的辐角主值的取值范围。

\Needspace{4\baselineskip}
\textbf{解析}

\Needspace{5\baselineskip}
令实根为 $b$。分离实、虚部得

\nopagebreak[4]
\[
b^2+4b+4=0,\qquad b+a=0,
\]

\Needspace{5\baselineskip}
所以 $b=-2$、$a=2$、$z=2-2i$。对 $c>0$，

\nopagebreak[4]
\[
\bar z(1-ci)=2(1+i)(1-ci)=2(1+c)+2(1-c)i.
\]

实部恒正。按固定扫描件采用的辐角区间 $[0,2\pi)$：当 $0<c\le1$ 时辐角从趋近但不取 $\pi/4$ 到 $0$；当 $c>1$ 时辐角在 $(7\pi/4,2\pi)$。合并为 $[0,\pi/4)\cup(7\pi/4,2\pi)$。

固定扫描件印刷页4-37的解答明确按 $[0,2\pi)$ 讨论辐角主值，并给出与此处一致的范围；卷内中间两句写作 $[0,\pi/2)$、$(3\pi/2,2\pi)$，属宽松表述，以最终答案为准。因此本题按扫描件所用约定作答；若改用常见的 $(-\pi,\pi]$ 约定，范围会写成 $(-\pi/4,\pi/4)$，不能据此把原卷结论静默改写。该独立推导本身仍待盲核。

\gkanswer{答案：$\left[0,\dfrac\pi4\right)\cup\left(\dfrac{7\pi}4,2\pi\right)$（卷内辐角区间取 $[0,2\pi)$）。}

\end{answer}
\begin{answer}
\gksolhead{20}{}

\textbf{题目}（本题满分16分）：

本题共2个小题，第1小题满分4分，第2小题满分12分。如图，在四棱锥 $P-ABCD$ 中，底面 $ABCD$ 是一直角梯形，$\angle BAD=90^\circ$，$AD\parallel BC$，$AB=BC=a$，$AD=2a$，且 $PA\perp$ 底面 $ABCD$，$PD$ 与底面成 $30^\circ$ 角。若 $AE\perp PD$，$E$ 为垂足。

（1）求证：$BE\perp PD$；

（2）求异面直线 $AE$ 与 $CD$ 所成角的大小（用反三角函数表示）。

原卷图（印刷页4-36）：
\nopagebreak[4]
\begin{center}% 1999 上海理科 Q20：依据固定扫描页的拓扑重绘；不表示长度/角度比例。
\begin{tikzpicture}[line join=round, line cap=round, scale=0.9]
  \tikzset{visible edge/.style={draw=black, line width=0.75pt}, hidden edge/.style={draw=black, dashed, line width=0.65pt}}
  \coordinate (P) at (0,3.05);
  \coordinate (A) at (0,0);
  \coordinate (B) at (-0.85,-1.65);
  \coordinate (C) at (1.45,-1.65);
  \coordinate (D) at (3.25,0.05);
  \coordinate (E) at (1.95,1.25);

  % Hidden base/back edges and the perpendicular projection to PD.
  \draw[hidden edge] (P)--(A);
  \draw[hidden edge] (A)--(B);
  \draw[hidden edge] (A)--(D);
  \draw[hidden edge] (A)--(E);
  % Visible pyramid edges and front/right base edges.
  \draw[visible edge] (P)--(B);
  \draw[visible edge] (P)--(C);
  \draw[visible edge] (P)--(D);
  \draw[visible edge] (B)--(C);
  \draw[visible edge] (C)--(D);

  \node[above] at (P) {$P$};
  \node[below right=4pt] at (A) {$A$};
  \node[left=3pt] at (B) {$B$};
  \node[below] at (C) {$C$};
  \node[right] at (D) {$D$};
  \node[above right] at (E) {$E$};
  \node[below=12mm] at (1.15,-1.65) {示意图，不按比例};
\end{tikzpicture}
\end{center}
\nopagebreak[4]
\Needspace{4\baselineskip}
\textbf{解析}

（1）$AB\perp AD$ 且 $PA\perp$底面，故 $AB$ 同时垂直于平面 $PAD$ 内相交的 $PA,AD$，从而 $AB\perp PD$。又 $AE\perp PD$，所以 $PD$ 垂直于平面 $ABE$ 内相交的 $AB,AE$，于是 $BE\perp PD$。

\Needspace{5\baselineskip}
（2）取 $a=1$，令 $A=(0,0,0)$、$B=(1,0,0)$、$D=(0,2,0)$、$C=(1,1,0)$、$P=(0,0,2/\sqrt3)$；最后一个坐标由斜线 $PD$ 与底面成 $30^\circ$ 得到。将 $A$ 正交投影到 $PD$ 得 $E=(0,1/2,\sqrt3/2)$。因此

\nopagebreak[4]
\[
\overrightarrow{AE}=(0,\tfrac12,\tfrac{\sqrt3}{2}),\quad
\overrightarrow{CD}=(-1,1,0),\quad
|AE|=1,\quad |CD|=\sqrt2,
\]

且 $\overrightarrow{AE}\cdot\overrightarrow{CD}=1/2$。两异面直线所成角取锐角，故其余弦为 $\sqrt2/4$。

\gkanswer{答案：（1）$BE\perp PD$；（2）$\arccos\dfrac{\sqrt2}{4}$。}

\end{answer}
\begin{answer}
\gksolhead{21}{}

\textbf{题目}（本题满分16分）：

本题共3个小题，第1小题满分6分，第2小题满分4分，第3小题满分6分。平地上有一条水沟，沟沿是两条长100米的平行线段，沟宽$AB$为2米，与沟沿垂直的平面与沟的交线是一段抛物线，抛物线顶点为$O$，对称轴与地面垂直，沟深1.5米，沟中水深1米。

（1）求水面宽；

（2）如图所示形状的几何体称为柱体，已知柱体的体积为底面积乘以高，问沟中的水有多少立方米？

（3）若要把这条水沟改挖（不准填土）成截面为等腰梯形的沟，使沟的底面与地面平行，则改挖后的沟底宽为多少米时，所挖的土最少？

原卷图（印刷页4-36）：
\nopagebreak[4]
\begin{center}% 1999 上海理科 Q21：保留水沟柱体与抛物线横截面的示意关系，不表示尺度。
\begin{tikzpicture}[line join=round, line cap=round, scale=0.82]
  \tikzset{visible edge/.style={draw=black, line width=0.75pt}, hidden edge/.style={draw=black, dashed, line width=0.65pt}}
  \coordinate (A)  at (0,0);
  \coordinate (B)  at (1.85,-0.55);
  \coordinate (O)  at (0.70,-1.75);
  \coordinate (A2) at (3.35,0.95);
  \coordinate (B2) at (5.20,0.35);
  \coordinate (O2) at (4.05,-1.10);
  \coordinate (M)  at (3.571,-0.088);

  % Near parabolic cross-section and prism length edges.
  \draw[visible edge] (A) .. controls (0.10,-0.90) and (0.33,-1.75) .. (O);
  \draw[visible edge] (O) .. controls (1.18,-1.75) and (1.68,-0.20) .. (B);
  \draw[visible edge] (A)--(A2);
  \draw[visible edge] (B)--(B2);
  \draw[visible edge] (O)--(O2);

  % The source shows the outer far profile, but its occluded inner segment is dashed.
  \draw[hidden edge] (A)--(B);
  \draw[visible edge] (A2) .. controls (3.425,0.615) and (3.495,0.245) .. (M);
  \draw[hidden edge] (M) .. controls (3.699,-0.647) and (3.843,-1.10) .. (O2);
  \draw[visible edge] (O2) .. controls (4.40,-1.10) and (4.90,-0.20) .. (B2);
  \draw[hidden edge] (A2)--(B2);

  \node[left] at (A) {$A$};
  \node[below=4pt] at (B) {$B$};
  \node[below] at (O) {$O$};
  \node[below=12mm] at (2.80,-1.75) {示意图，不按比例};
\end{tikzpicture}
\end{center}
\nopagebreak[4]
\Needspace{4\baselineskip}
\textbf{解析}

设抛物线横截面为 $y=\frac32x^2$，顶点为沟底 $O$，地面高 $y=3/2$，沟沿对应 $x=\pm1$。

（1）水深为1，水面 $y=1$ 与抛物线交于 $x=\pm\sqrt{2/3}$，故水面宽为 $2\sqrt{2/3}=2\sqrt6/3$ 米。

\Needspace{5\baselineskip}
（2）横截面水面积为

\nopagebreak[4]
\[
2\int_0^{\sqrt{2/3}}\left(1-\frac32x^2\right)\,dx=\frac{4\sqrt6}{9}\text{\text{平方米}}.
\]

水沟长100米，故水量为 $100\cdot4\sqrt6/9=400\sqrt6/9$ 立方米。

\Needspace{5\baselineskip}
（3）最优梯形的底边、顶边分别可取 $y=0$ 与地面 $y=3/2$：若候选底边低于 $y=0$，将其上移并保留它在 $y=0$ 处的宽度；若顶边高于地面，则截去地面以上部分。两种调整都仍包含原沟且面积更小。记右侧底、顶半宽分别为 $a,c$，并令归一化高度 $s=2y/3\in[0,1]$。右侧直边在该高度的半宽是

\nopagebreak[4]
\[
h(s)=a+ds,\qquad d=c-a.
\]

\Needspace{5\baselineskip}
原抛物线在高度 $y=3s/2$ 的半宽为 $\sqrt{s}$。令 $x=\sqrt{s}$，不填土的充要条件是

\nopagebreak[4]
\[
h(s)\ge\sqrt{s}\ (0\le s\le1)\quad\Longleftrightarrow\quad q(x)=a+dx^2-x\ge0\ (0\le x\le1).
\]

这就是“不准填土”的逐高度包含条件，保证每一层原沟均被新沟覆盖而不需填土。

梯形总截面积为 $S=\frac32(a+c)=\frac32(2a+d)$。对所有可行斜率分段求下界：

\begin{itemize}
\item 若 $d\le0$，则 $q'(x)=2dx-1<0$，由 $q(1)=c-1\ge0$ 得 $c\ge1$，所以 $a+c=2c-d\ge2$。
\item 若 $0<d\le1/2$，$q$ 在 $[0,1]$ 上递减，仍须 $c\ge1$；于是 $a\ge1-d$，从而 $a+c=2a+d\ge2-d\ge3/2$。
\item 若 $d\ge1/2$，$q$ 的顶点 $x=1/(2d)$ 位于 $[0,1]$，可行性要求 $a\ge1/(4d)$。故
\end{itemize}

\nopagebreak[4]
\[
a+c=2a+d\ge d+\frac1{2d}\ge\sqrt2,
\]

\Needspace{5\baselineskip}
最后一步由均值不等式成立，且仅在 $d=1/\sqrt2$ 取等号。因为 $\sqrt2<3/2<2$，全局最小值出现在此分支；取 $a=1/(4d)=\sqrt2/4$、$c=a+d=3\sqrt2/4$ 可实现该下界。令 $t=1/\sqrt2$，便有 $a=t/2,d=1/(2t)$。此时

\nopagebreak[4]
\[
h(s)-\sqrt{s}=\frac{s+t^2}{2t}-\sqrt{s}=\frac{(\sqrt{s}-t)^2}{2t}\ge0,\qquad t=1/\sqrt2,
\]

亦直接验证新沟逐高度包含原沟。因此全局最优沟底宽 $2a=\sqrt2/2$ 米，上口半宽为 $3\sqrt2/4$ 米，完整横截面积为 $3\sqrt2/2$ 平方米。

\Needspace{5\baselineskip}
原题没有要求新沟口仍固定为2米。若另加“沟口宽度固定为2米”的约束，则 $c=1$、$d=1-a$。当 $a<1/2$ 时 $d>1/2$，顶点条件要求 $a\ge1/(4d)$，即 $4d(1-d)\ge1$，这与 $d>1/2$ 矛盾；因此 $a\ge1/2$。取 $a=1/2,d=1/2$ 时 $q(x)=\frac12(x-1)^2\ge0$，达到该约束下最小底宽1米，完整横截面积为 $9/4$ 平方米。原沟截面积为

\nopagebreak[4]
\[
\int_{-1}^{1}\left(\frac32-\frac32x^2\right)\,dx=2\text{\text{平方米}}.
\]

因此自由沟口与固定2米沟口的新增挖方截面积分别为 $3\sqrt2/2-2$ 平方米和 $1/4$ 平方米；前者严格较小（总面积平方比较为 $9/2<81/16$）。固定口宽的1米不是自由沟口解释下的答案。按本库独立复核（逐高度包含条件与全局最小证明）：自由沟口解释下的最小底宽为 $\sqrt2/2$ 米；若新沟口固定 2 米，则为 $1$ 米。来源未明示改挖后的沟口尺寸，题意仍有来源歧义，本推导为编者独立推导，不等同官方答案。

\gkanswer{答案：（1）$\dfrac{2\sqrt6}{3}$ 米；（2）$\dfrac{400\sqrt6}{9}$ 立方米；（3）$\dfrac{\sqrt2}{2}$ 米。}

\end{answer}
\begin{answer}
\gksolhead{22}{}

\textbf{题目}（本题满分18分）：

本题共3个小题，第1小题满分6分，第2小题满分6分，第3小题满分6分。设椭圆 $C_1$ 的方程为 $\frac{x^2}{a^2}+\frac{y^2}{b^2}=1\ (a>b>0)$，曲线 $C_2$ 的方程为 $y=\frac1x$，且 $C_1$ 与 $C_2$ 在第一象限内只有一个公共点 $P$。

（1）试用 $a$ 表示点 $P$ 的坐标；

（2）设 $A,B$ 是椭圆 $C_1$ 的两个焦点，当 $a$ 变化时，求 $\triangle ABP$ 的面积函数 $S(a)$ 的值域；

（3）记 $\min\{y_1,y_2,\ldots,y_n\}$ 为 $y_1,y_2,\ldots,y_n$ 中最小的一个。设 $g(a)$ 是以椭圆 $C_1$ 的半焦距为边长的正方形的面积，试求函数 $f(a)=\min\{g(a),S(a)\}$ 的表达式。

\Needspace{4\baselineskip}
\textbf{解析}

\Needspace{5\baselineskip}
（1）代入 $y=1/x$ 后令 $u=x^2>0$，交点满足

\nopagebreak[4]
\[
\frac{u}{a^2}+\frac{1}{ub^2}=1
\iff b^2u^2-a^2b^2u+a^2=0.
\]

\Needspace{5\baselineskip}
第一象限内恰有一个交点，二次方程须有重根，判别式为零给出 $a^2b^2=4$，即 $ab=2$。重根为 $u=a^2/2$，又 $a>b>0$，故 $a>\sqrt2$ 且

\nopagebreak[4]
\[
P=\left(\frac a{\sqrt2},\frac{\sqrt2}{a}\right).
\]

\Needspace{5\baselineskip}
（2）椭圆半焦距 $c=\sqrt{a^2-b^2}=\sqrt{a^2-4/a^2}$。以 $AB=2c$ 为底、点 $P$ 的纵坐标为高，

\nopagebreak[4]
\[
S(a)=c\frac{\sqrt2}{a}=\sqrt{2\left(1-\frac4{a^4}\right)}.
\]

当 $a$ 从 $\sqrt2$ 右侧趋近时 $S(a)\to0$，当 $a\to+\infty$ 时 $S(a)\to\sqrt2$，且函数严格递增，因此值域为 $(0,\sqrt2)$。

\Needspace{5\baselineskip}
（3）$g(a)=c^2=a^2-4/a^2$。在 $a>\sqrt2$ 上两式均正；解 $g(a)=S(a)$ 并平方，令 $u=a^4$ 得

\nopagebreak[4]
\[
u^2-10u+24=0\iff (u-4)(u-6)=0.
\]

\Needspace{5\baselineskip}
因 $u>4$，唯一相关分界点是 $a=\sqrt[4]6$。在 $\sqrt2<a\le\sqrt[4]6$ 时 $g(a)\le S(a)$，其后 $S(a)<g(a)$，故

\nopagebreak[4]
\[
f(a)=\begin{cases}
a^2-\dfrac4{a^2},&\sqrt2<a\le\sqrt[4]6,\\[4pt]
\sqrt{2\left(1-\dfrac4{a^4}\right)},&a>\sqrt[4]6.
\end{cases}
\]

\gkanswer{答案：（1）$\left(\dfrac{a}{\sqrt2},\dfrac{\sqrt2}{a}\right)$（$a>\sqrt2$）；（2）$S(a)$ 的值域为 $(0,\sqrt2)$；（3）$\displaystyle f(a)=\min\{g(a),S(a)\}=\begin{cases}a^2-\dfrac{4}{a^2},&\sqrt2<a\le\sqrt[4]{6},\\[2pt]\sqrt{2\left(1-\dfrac{4}{a^4}\right)},&a>\sqrt[4]{6}.\end{cases}$}

\end{answer}